\documentclass[11pt]{article}
\usepackage[margin=1in]{geometry}
\usepackage{booktabs}
\usepackage{microtype}
\usepackage[hyphens]{url}
\sloppy

\title{FastBandSolver: review preparation}
\date{}

\begin{document}
\maketitle

Questions a reviewer is likely to ask, with prepared answers. Numbers are from the paper's measurements (Apple M4
Max, the run after the row-single fix) and tdoku's published results for an Intel Core i7-1065G7.

\section*{Is most of the benefit just running the tdoku algorithm on an M4 Max?}

No. There are two separate effects, and on the hard puzzles they matter about equally.

First, FastBandSolver does not run tdoku's algorithm. tdoku uses 256-bit SIMD with triads and message passing;
FastBandSolver is JCZSolve's scalar band engine plus three changes (peer-count branching, the stack rule and a
fewest-candidates fallback). tdoku itself gets slower on the M4, about 0.88 times its published Intel speed on the hard
sets.

On forum hardest 1106, with each solver's speed relative to tdoku on the same machine:

\begin{center}
\begin{tabular}{lr}
\toprule
 & speed relative to tdoku \\
\midrule
JCZSolve on Intel (published) & 0.41$\times$ \\
JCZSolve on M4 (measured) & 0.68$\times$ \\
our engine with the three changes switched off, on M4 & 0.62$\times$ \\
FastBandSolver on M4 & \textbf{1.17$\times$} \\
FastBandSolver on Intel (projected) & 0.70$\times$ \\
\bottomrule
\end{tabular}
\end{center}

\begin{itemize}
\item \textbf{The platform effect, about 1.7$\times$.} Moving from Intel to the M4 lifts the scalar design from 0.41 to
  0.68 of tdoku's speed, because scalar code gains on the M4 and tdoku does not.
\item \textbf{The algorithmic changes, about 1.7--1.9$\times$.} Peer-count branching, the stack rule and the
  fewest-candidates fallback take it from 0.62--0.68 to 1.17.
\item \textbf{Neither alone is enough.} JCZSolve on the M4 still loses to tdoku (0.68$\times$), and the changes
  projected onto Intel still lose (0.70$\times$). Both are needed.
\end{itemize}

On the other data sets:

\begin{center}
\begin{tabular}{lrr}
\toprule
data set & platform shift & algorithm gain (vs JCZSolve on M4) \\
\midrule
magictour top 1465 & 1.56$\times$ & 1.45$\times$ \\
forum hardest 1905 11+ & 1.67$\times$ & 1.40$\times$ \\
17 clue & 1.62$\times$ & 1.33$\times$ \\
kaggle & 1.63$\times$ & 1.36$\times$ \\
\bottomrule
\end{tabular}
\end{center}

On 17-clue the platform alone is enough: JCZSolve already beats tdoku there on the M4 (1.22$\times$), so the
algorithmic changes add margin rather than create the win.

\paragraph{Caveat.} The platform figure is the least robust. The Intel numbers are tdoku's published ones, from a
different machine and benchmark program, each the best of over a dozen compiler configurations, so ``about
1.7$\times$'' is a cross-machine estimate. The algorithmic figure was measured on the same machine with interleaved
runs.

\paragraph{Summary for the paper.} The M4 closes most of the gap between scalar band solvers and tdoku, and the
algorithmic changes close the rest and pull ahead. The paper states this across Sections 6.1 and 7; it could be
stated in one paragraph in the discussion.


\section*{Has anybody done this before?}

We searched GitHub, the Sudoku Players' Forum (forum.enjoysudoku.com, down during the search, so read through
incomplete Wayback Machine snapshots), arXiv and the constraint-programming literature, in September 2026. Two 2026
projects are close relatives; neither does what we do, but both must be cited.

\paragraph{Solvers claimed to be faster than tdoku.}
\begin{itemize}
\item \textbf{fastdoku} (rocketbombs, August 2026, \url{github.com/rocketbombs/fastdoku}). A Rust solver that combines
  a JCZSolve-style band engine with a port of tdoku's triad engine, choosing per puzzle. It claims 1.06--1.63$\times$
  over the faster of tdoku and rust\_sudoku on all eight of tdoku's corpora, on x86. Its benchmark is \emph{time to
  first solution} (limit 1) on both sides, so it is consistent but not comparable to our limit-2 results. Its README
  says its lead is x86-specific; on ARM it has a NEON port of its triad engine but compares only against rust\_sudoku,
  stating that tdoku cannot be built on ARM. It is the closest prior work.
\item \textbf{RKCR7 / autoresearch-sudoku} (Rkcr7, March 2026,
  \url{github.com/Rkcr7/autoresearch-sudoku}). A Rust
  JCZSolve-family solver built by an AI coding agent, claiming to beat tdoku on 4 of 6 data sets on an AMD Ryzen with
  AVX-512. Its \texttt{solve()} stops at the first solution, while tdoku was run through tdoku's own benchmark program,
  which checks uniqueness (limit 2) by default. We confirmed this from its log: tdoku made 61.67 guesses per puzzle on
  forum hardest 1905 11+, exactly tdoku's count at limit 2 on our machine (34.17 at limit 1). So its comparison is
  first solution against uniqueness check, which favors it; the README does not say so.
\item \textbf{gz\_sudoku} (shines77, 2021--2024): a C++ rewrite of JCZSolve; its README concedes tdoku is fastest on
  hard puzzles.
\end{itemize}
We found no claim of beating tdoku with a uniqueness check (limit 2), on any platform.

\paragraph{tdoku on ARM.} Besides the SIMDe-based fork we already cite, fastdoku ports tdoku's \emph{design} to NEON
in Rust (not tdoku's code), and a Rust port (brki/rust-tdoku) runs on ARM with scalar fallbacks only. We found no
native NEON build of tdoku itself other than ours, and no benchmark of tdoku, JCZSolve or fsss2 on bare-metal Apple
Silicon or other ARM servers.

\paragraph{Our three changes.}
\begin{itemize}
\item \textbf{Peer-count branching} (the two-candidate cell with the most unsolved peers): this is the classic
  \emph{degree} tie-break for minimum remaining values (Br\'elaz 1979), and it has been evaluated on Sudoku in generic
  constraint solvers (a 2015 Groningen bachelor's thesis by Bakker; an arXiv paper, DiBS, 2026). No precedent found in
  a fast solver. Every band solver we read (JCZSolve, rust\_sudoku, RKCR7, fastdoku) branches on the first two-candidate
  cell; fsss2 has related experimental strategies, off by default. Degree-style heuristics are standard in constraint
  programming in general.
\item \textbf{Stack rule} (bands matched to columns within each stack): the constraint itself is in Simonis,
  ``Sudoku as a Constraint Problem'' (2005), Section 5.3, which gives the rows--blocks matching (our band rule) and
  notes the columns--blocks version (our stack rule), solved with a general matching algorithm. We read the paper to
  confirm. What is new is applying it by table lookup in a scalar band solver; tdoku reaches the same deductions
  through its vertical-band configurations.
\item \textbf{Fewest-candidates fallback}: near precedent. RKCR7 and fastdoku take the fewest-candidate cell among the
  first unsolved cell of each band, rust\_sudoku checks up to three cells, and the rule itself is classic
  minimum-remaining-values.
\end{itemize}
The \emph{band} rule (rows matched to boxes within a band) is not new: fastdoku presents its exact version as a
strengthening over JCZSolve, but we checked JCZSolve's table (\texttt{TblComplexMask}) against the perfect-matching
definition on all 512 patterns and it is identical, so JCZSolve already had it. The single lookup table for row singles
also appears to be new; rust\_sudoku and RKCR7 combine two tables as JCZSolve does.

\paragraph{Academic work.} A second search covered arXiv (all 54 results for ``sudoku solver''), the CP and SAT
literature and theses; Semantic Scholar, DBLP and some publishers' pages were blocked. No paper proposes a fast
9$\times$9 solver or benchmarks against tdoku, JCZSolve or fsss2; papers that mention tdoku use it only to rate puzzle
difficulty (for example the Hierarchical Reasoning Model paper, 2025). Eppstein (FUN 2012) solves single-digit
subproblems faster than the pattern overlay method, which supports our template negative result, and Simonis's
\emph{shaving} is the lookahead we found too costly. We found no study comparing Sudoku solvers across processor
families, and none on branching heuristics under a uniqueness requirement.

\paragraph{What we can claim.} The first result, as far as we can find, of a solver beating tdoku with a uniqueness
check (limit 2), measured against tdoku's own code on the same machine; the first use of Simonis's stack constraint and
of the degree tie-break in a fast band solver, with measurements of each; and the first native NEON build of tdoku. The
rules and the heuristic themselves are not ours. We cannot claim to be first to beat tdoku
under any setting (fastdoku does, at limit 1 on x86), first with an exact band rule (JCZSolve), or first with a
fewest-candidates fallback.

\paragraph{Limits of the search.} The forum was offline and its archive incomplete after 2019, so forum-only work on
guess strategies may exist that we could not see, and closed-source solvers could not be checked.

\paragraph{Action for the paper.} Cite fastdoku and RKCR7 in the related work, with the limit-1 caveat for both and the
mismatched baseline for RKCR7.

\end{document}
